% Contact-establishment reference generation, seen along a surface tangent with
% the tool face parallel to the surface. The face is flat here, so the selected
% tool point is the face centre and lies on the end-effector axis below the TCP.
% The commanded point advances past the surface, so the endpoint is a spring
% reference rather than an expected travel.
% The two offsets drawn are the same rigid offset: p_TCP(t_start) above
% p_Tool(t_start), and p_d(t_end) above p_Tool,d(t_end). Drawing both makes the construction
% of Equation (3.20) visible, so neither segment carries a label: the product
% R_EE(t_start) r_Tool,EE was the longest item in the picture and Equation
% (3.14) states the relation on the same page. The commanded travel carries no
% label either: the arrow runs from p_Tool,0 to p_Tool,d(t), which is what
% Equation (3.18) states, and s_c is the distance between those two points.
% Drawn in black throughout. The chain is read from the line styles instead of
% from colour: the surface and the tool face are the heaviest strokes, the rigid
% tool offset is a plain line, the projection onto the surface is dashed because
% it is construction, and the commanded travel is the one arrow on the chain.
% Labels alternate left and right of the axis so that no two sit on one side at
% the same height.
% Uses only tikz + arrows.meta, which config/packages.tex loads.
\begin{tikzpicture}[
    scale=1.55,
    font=\footnotesize,
    tool/.style={draw=black, line width=1.0pt},
    surf/.style={draw=black, line width=1.0pt},
    off/.style={draw=black, line width=0.8pt},
    proj/.style={draw=black, line width=0.7pt, densely dashed},
    cmd/.style={-{Latex[length=2.0mm]}, black, line width=1.0pt},
    nrm/.style={-{Latex[length=2.0mm]}, black, line width=0.8pt},
    leader/.style={draw=black, thin},
    pt/.style={fill=black, draw=black},
    lbl/.style={font=\footnotesize, inner sep=1.8pt},
  ]

  % ---- surface
  \draw[surf] (0.40,0) -- (4.75,0);
  \node[lbl, anchor=west] at (4.82,0) {Surface};

  % ---- the surface normal, rising from the surface clear of the chain
  \draw[nrm] (0.72,0) -- (0.72,0.68);
  \node[lbl, anchor=west] at (0.81,0.40) {$n_s$};

  % ---- tool face at the clearance transition, parallel to the surface
  \draw[tool] (1.50,1.05) rectangle (3.50,1.50);

  % ---- TCP, and the rigid offset down to the selected tool point
  \draw[off] (2.50,2.15) -- (2.50,1.05);
  \node[pt] at (2.50,2.15) [circle, inner sep=1.6pt] {};
  \node[lbl, anchor=west] at (2.66,2.15) {$p_{\mathrm{TCP}}(t_{\mathrm{start}})$};

  % ---- the selected tool point is the face centre while the face is flat
  \node[pt] at (2.50,1.05) [circle, inner sep=1.7pt] {};
  \draw[leader] (2.44,1.00) -- (2.22,0.84);
  \node[lbl, anchor=east] at (2.18,0.82) {$p_{\mathrm{Tool}}(t_{\mathrm{start}})$};

  % ---- its projection onto the configured surface reference
  \draw[proj] (2.50,1.05) -- (2.50,0);
  \node[pt] at (2.50,0) [circle, inner sep=1.7pt] {};
  \node[lbl, anchor=west] at (2.64,0.24) {$p_{\mathrm{Tool},0}$};

  % ---- commanded motion of that point, continuing past the surface
  \draw[cmd] (2.50,0) -- (2.50,-1.45);
  \node[pt] at (2.50,-1.45) [circle, inner sep=1.7pt] {};
  \node[lbl, anchor=west] at (2.64,-1.45) {$p_{\mathrm{Tool},d}(t_{\mathrm{end}})$};

  % ---- the TCP target that realises it, one rigid offset above
  \node[pt] at (2.50,-0.40) [circle, inner sep=1.6pt] {};
  \node[lbl, anchor=east] at (2.34,-0.40) {$p_d(t_{\mathrm{end}})$};

\end{tikzpicture}
